\chapter{\IfLanguageName{dutch}{Stand van zaken}{State of the art}}%
\label{ch:stand-van-zaken}

% Tip: Begin elk hoofdstuk met een paragraaf inleiding die beschrijft hoe
% dit hoofdstuk past binnen het geheel van de bachelorproef. Geef in het
% bijzonder aan wat de link is met het vorige en volgende hoofdstuk.

% Pas na deze inleidende paragraaf komt de eerste sectiehoofding.

\section{Literatuurstudie}%
\label{sec:literatuurstudie1}
\subsection{Context}
Het 24flow platform is gemaakt in het Salesforce ecosysteem voor make-to-order bedrijven (\autocite{Nelfiyanti2025}), waar workflows vaak complex zijn en variëren per order. Dit betekent dat een fabriek niet precies weet welk product hij zal maken, tot een gebruiker effectief iets koopt of een order plaatst. Een gewone fabriek maakt bijvoorbeeld 1000 dezelfde stoelen, een make-to-order bedrijf zou dan tien gewone stoelen kunnen maken, één groene, één met extra poten enz. Er is met andere woorden geen standaard pad dat de productie neemt. Hierdoor is er dan ook een noodzaak voor digitale oplossingen die heel specifiek voor verschillende productieprocessen kunnen worden gebruikt en het mogelijk maken om elk uniek order efficiënt te beheren. 
\\

24flow is zo een oplossing, het maakt gebruik van de flexibiliteit van Salesforce om voor elk individueel order een dynamisch productiepad te berekenen. Het verbindt engineering, planning en de werkvloer in één visuele en traceerbare workflow. Het zorgt ervoor dat operatoren en teamleden over de juiste informatie en tools beschikken om efficiënter te werken. Hierdoor krijgt elk uniek order, ondanks de mogelijke variatie aan productiestappen telkens de juiste informatie en de meest efficiënte route doorheen de fabriek.
\\

Het platform is gebouwd binnen het Salesforce ecosysteem met behulp van Apex, de programmeertaal van Salesforce zelf. Salesforce is een algemeen erkend customer relationship management system (CRM)-oplossing (\autocite{Pookandy2024}). Een CRM-systeem functioneert als een centraal platform voor het beheren van alle data tussen een bedrijf en zijn klanten. Salesforce voorziet daarnaast allerlei tools om op een flexibele manier data, workflows, gebruikers... te beheren en allerlei opties te geven aan bedrijven om heel specifieke oplossingen te kunnen implementeren die verder gaan dan enkel klantbeheer. 24flow benut deze flexibiliteit om complexe productiegegevens om te zetten in bruikbare informatie voor de werkvloer. Dit betekent bijvoorbeeld dat een operator of werknemer op de werkvloer enkel de instructies en data ziet die relevant zijn voor het specifieke product dat op dat moment bewerkt wordt, volledig afgestemd op de huidige status van het productieproces.


\subsection{Onboarding}
Onboarding is het proces waarbij een nieuwe gebruiker voor het eerst in contact komt met een applicatie, of andere toepassing en de manier waarop die gebruiker hierbij wordt begeleid. Die begeleiding of sturing kan verschillende vormen aannemen, het kan gebeuren aan de hand van een menselijke uitleg, walkthroughs of handleidingen. \autocite{Madhwacharyula2023} 

Hierbij zijn er telkens verschillende voor- en nadelen te vinden:

\begin{itemize}
	\item \textbf{Menselijke uitleg:} Een menselijke uitleg is zeer interactief en zorgt voor de meeste flexibiliteit, er kunnen complexe vragen gesteld worden en er kan heel diepgaand begeleid worden. Wel is het heel tijdsintensief en niet schaalbaar voor grotere groepen.
	\item \textbf{Walkthroughs:} Dit zijn stappenplannen die gebruikers stapsgewijs door een applicatie of interface leiden. Deze hulp is direct aanwezig in de applicatie zelf, dit kan bijvoorbeeld gebruikt worden voor het opzetten van een nieuwe configuratie of applicatie. Als de gebruiker daarentegen vastloopt of vragen heeft, is hier weinig ruimte toe, het is veelal een zelfsturend proces.
	\item \textbf{Handleiding:} Dit is het meest statische hier, het is bedoeld als zelfsturend werk, als document om zonder enige vorm van interactie een gebruiker te helpen met een taak. Vaak is dit ook de manier om een uitleg nog eens te kunnen raadplegen, een naslagwerk dat altijd beschikbaar is om terug iets in op te zoeken.
	\item \textbf{Video tutorial:} Een handleiding kan echter ook de vorm van een video aannemen. Hier heb je een soort combinatie van walkthrough, menselijke uitleg en handleiding. Stapsgewijs wordt er iets uitgelegd en de gebruiker kan op zijn eigen tempo werken, hij kan immers pauzeren wanneer nodig. Wel is hier geen interactie en ziet een tutorial meestal de specifieke situatie van de gebruiker niet.
\end{itemize}

Onboarding speelt een grote rol bij de first time user experience (FTUE), de eerste interactie die een gebruiker heeft met een bepaalde applicatie, en waar de mening van die gebruiker over de applicatie fundamenteel gevormd wordt. 
\autocite{CascaesCardoso2017} Een goed ontworpen onboardingproces helpt gebruikers niet enkel om de basisfunctionaliteiten te begrijpen, maar verhoogt de kans dat zij de applicatie blijven gebruiken.
Effectieve onboarding betekent immers beter begrip, minder frustraties en een goede gebruikerservaring van de applicatie op de lange termijn.
Onderzoek van \textcite{ChukwunonsoNweke2025} over verschillende onboarding technieken toont dat walkthroughs die een gebruiker stap per stap begeleiden ervoor zorgen dat die gebruiker de nieuwe kennis beter begrijpt en onthoudt dan bij een eenmalige menselijke uitleg.
\\

Goede onboarding zorgt ervoor dat alle onderdelen van een applicatie stapsgewijs geïntroduceerd worden, en telkens met de nodige context getoond worden. Dit betekent dat een gebruiker de juiste functies en informatie krijgt die op dat moment relevant zijn voor hem, bijvoorbeeld gekoppeld aan de taak die hij op dat moment aan het uitvoeren is. Verder kunnen tutorials en walkthroughs hierbij een grote hulp zijn om gebruikers gemakkelijk wegwijs te maken met een bepaalde interface of applicatie.
\autocite{Kedzielska2023}
\autocite{Melenciano2025}
\\

In tegenstelling hiertoe kan slechte onboarding  leiden tot frustratie, verwarring en zelfs het verlaten van een platform, dus het verlies van een klant.
\autocite{CascaesCardoso2017}
\\

Moderne onboarding gebeurt tegenwoordig vaak via in-app systemen, contextgevoelige prompts, interactieve tutorials en walkthroughs. Deze aanpak zorgt ervoor dat de gebruiker op het juiste moment de juiste informatie of hulp krijgt. \autocite{Acar2020, Akiki2019}

Deze methoden zijn effectief omdat ze een gebruiker minder overweldigen met alle informatie op hetzelfde moment:
\begin{itemize}
	\item \textbf{In-app systemen:} Het gebruik van in-app systemen voor onboarding betekent dat de nodige hulp direct in de interface is geïntegreerd.
	\item \textbf{Contextgevoelige prompts:} Hierbij spelen contextgevoelige prompts dan in op de specifieke acties die een gebruiker uitvoert, dit kan bijvoorbeeld als een pop-up zijn wanneer een gebruiker een ingewikkelde functie opent.
	\item \textbf{Interactieve tutorials en walkthroughs:} Deze combineren deze systemen: zij laten de gebruiker stapsgewijs handelingen uitvoeren en geven directe feedback binnen de applicatie met behulp van onder andere zo'n prompts. Ze zijn in staat om de uitgevoerde handelingen van de gebruiker te controleren en direct bij te sturen waar nodig.
\end{itemize}

\subsection{In-App User Guidance}
Het concept in-app user guidance, of letterlijk het sturen van gebruikers tijdens het gebruik van een applicatie, betekent dat een gebruiker alle nodige info en hulp kan vinden binnen de applicatie. Volgens \textcite{Acar2020} is het voor gebruikers moeilijk en demotiverend indien hulp enkel beschikbaar is in bijvoorbeeld handleidingen, tutorials... en gebruikers zelf actief op zoek moeten gaan naar de juiste oplossing. Ze duiden erop dat gebruikers hulp zouden moeten krijgen op het exacte moment dat ze die nodig hebben en dat deze hulp dus context specifiek moet aangeboden worden. Bijvoorbeeld wanneer een gebruiker voor het eerst een nieuw onderdeel van de applicatie gebruikt, krijgt hij een korte tip of aanwijzing over dat specifieke onderdeel. Hierdoor krijgt de gebruiker dus relevante informatie precies op het juiste moment, zonder dat hij zelf overdreven ver moet zoeken.
\\

\textcite{Akiki2019} dan heeft het over de praktische implementatie van zo'n user guidance, de studie toont aan dat hulp voor gebruikers direct kan ingebouwd worden in de gebruikersinterface via context gebonden prompts en interactieve elementen. Dit onderzoek laat eveneens zien dat gebruikers die op deze manier geholpen worden, efficiënter kunnen werken en minder vastlopen dan wanneer ze zelf op zoek moeten naar alle juiste oplossingen.
\\

Daarbovenop toont onderzoek over digital adoption platforms (DAP) dat het integreren van contextspecifieke en interactieve ondersteuning binnen enterprise applicaties het aanleren van zo'n complexe systemen verbetert.\autocite{Handrich2024}
\\

Een DAP is een softwarelaag bovenop een applicatie die de gebruiker actief stuurt binnen deze applicatie. Het begeleidt een gebruiker bijvoorbeeld door de juiste knoppen te laten oplichten, of door pop-ups te tonen op het juiste moment. Het maakt dus gebruik van allerlei in-app guidance functies, walkthroughs en dergelijke meer om een gebruiker wegwijs te maken in vaak complexe systemen. Hoewel een DAP vaak een externe tool is heeft Salesforce deze functionaliteit rechtstreeks in zijn systeem geïntegreerd.
\\

Immers kunnen de bestaande tools van Salesforce voor in-app guidance, zoals walkthroughs, interactieve prompts enz. gezien worden als een dergelijke DAP laag die gebruikers in staat stelt om actief begeleid te worden terwijl ze in de applicatie werken.
\clearpage

\section{Salesforce Analyse}%
\label{sec:Salesforce analyse}
Voordat er een technische oplossing voor de onboarding van 24flow kan worden uitgewerkt, is het essentieel om te begrijpen welke functionaliteiten Salesforce standaard aanbiedt. Eerst worden de kerncomponenten van de Salesforce In-App Guidance belicht, zoals prompts en walkthroughs en de configuratie hiervan. Vervolgens wordt besproken welke overzicht en analyse tools Salesforce biedt om guidance te analyseren en beheren. 
\\

\subsection{Technologisch Kader en Begrippen}
Om de technische context van dit onderzoek te verduidelijken, worden enkele belangrijke componenten en begrippen van het Salesforce-ecosysteem hieronder kort toegelicht.

\begin{itemize}
	\item \textbf{Lightning Experience (LEX):} De moderne, component-gebaseerde gebruikersinterface van Salesforce. LEX fungeert als een modulaire werkomgeving waarin componenten op flexibele pagina-layouts kunnen worden geplaatst.
	
	\item \textbf{Lightning Web Components (LWC):} Het moderne framework voor het bouwen van gebruikersinterfaces binnen Salesforce, gebaseerd op de huidige webstandaarden (HTML, moderne JavaScript en CSS).
	
	\item \textbf{Salesforce Lightning Design System (SLDS):} Een uitgebreid framework van CSS-richtlijnen en UI-componenten. SLDS zorgt voor een consistente visuele stijl en gebruikerservaring die naadloos aansluit op de standaard Salesforce-interface.
	
	\item \textbf{Apex:} Een objectgeoriënteerde, sterk getypeerde programmeertaal die op de servers van Salesforce wordt uitgevoerd. Het stelt ontwikkelaars in staat om complexe logica te implementeren.
	
	\item \textbf{Custom Metadata Types (MDT):} Een specifiek type record dat wordt gebruikt voor het opslaan van configuratiedata in plaats van transactionele bedrijfsdata. Metadata is uniek omdat de records zelf als onderdeel van de applicatiestructuur worden beschouwd.
	
\end{itemize}

\subsection{In-App Guidance Tools}

Volgens de officiële specificaties van Salesforce is In-App Guidance functionaliteit opgebouwd uit twee fundamentele hoofdelementen: \textit{Prompts} en \textit{Walkthroughs}, die ondersteund worden door uitgebreide configuratiemogelijkheden.

\subsubsection*{1. Prompts}
Een prompt is een losstaande aanwijzing die is ontworpen om de aandacht van de gebruiker te vestigen op specifieke elementen, acties of updates binnen het systeem. Salesforce onderscheidt drie specifieke prompts, elk met een eigen functioneel doel:

\begin{itemize}
	\item \textbf{Targeted Prompts:} Dit zijn vensters die visueel verankerd zijn aan een specifiek UI-element, zoals een invulveld, knop of tabblad. De prompt wijst met een pijl direct naar het bijpassende element en maakt de rest van de pagina grijs, hierbij zijn er twaalf verschillende posities die de prompt kan innemen ten opzichte van het element. 
	\\
	
	\textit{Toepassing:} Dit is uiterst effectief om een gebruiker heel gericht te sturen naar bepaalde knoppen of velden die op dat moment essentieel zijn. (bijv. een prompt die naar een nieuw toegevoegde knop wijst, waardoor de gebruiker direct op de hoogte is van de nieuwe functionaliteit.  )
	
	\item \textbf{Floating Prompts:} Dit zijn pop-upvensters die niet aan een specifieke knop zijn gekoppeld, maar los over de interface zweven. Deze prompt kan op negen verschillende voorgedefinieerde locaties op een pagina worden weergegeven. Dit soort venster neemt weinig ruimte in en verdwijnt onmiddellijk wanneer een gebruiker naar een andere pagina navigeert of de prompt weg klikt.
	\\
	
	\textit{Toepassing:} Deze worden doorgaans ingezet als kleinere tooltips of informatie over bepaalde functionaliteiten.(bijv. een werknemer opent een lange lijst met orders en een floating prompt verschijnt met een korte uitleg over de mogelijkheid om filters toe te passen op lijsten.)
	
	\item \textbf{Docked Prompts:} Een docked prompt verschijnt standaard geminimaliseerd of geopend in de rechterbenedenhoek van het scherm. Het unieke aan dit type is dat het op het scherm blijft staan terwijl de gebruiker de rest van de applicatie bedient. De gebruiker kan het venster openklappen of sluiten naar behoefte.
	\\
	
	\textit{Toepassing:} Omdat dit type meer tekst en zelfs media ondersteunt, is het ideaal voor uitgebreide handleidingen, checklists of zelfs het gebruiken van video-tutorials. (bijv. een video met de voornaamste laatste nieuwe features van de voorbije release die beschikbaar zijn.)
\end{itemize}
\begin{table}[h!]
	\centering
	\caption{Functionele vergelijking van Salesforce In-App Guidance prompt-types}
	\label{tab:prompt_vergelijking}
	\small
	\begin{tabular}{@{}lccc@{}}
		\toprule
		\textbf{Kenmerk } & \textbf{Floating Prompt} & \textbf{Targeted Prompt} & \textbf{Docked Prompt} \\ \midrule
		Rich Text Support            & Basis     & Basis     & Uitgebreid* \\
		Afbeeldingen ondersteund     & Ja        & Ja        & Ja           \\
		Embedded Video               & Nee       & Nee       & Ja           \\
		Action / Dismiss Button      & Ja        & Ja        & Ja           \\
		Snooze-icoon                 & Ja        & Ja        & Ja           \\ 
		Gedrag bij navigatie         & Verdwijnt & Verdwijnt & Blijft staan \\
		Grootte op pagina            & Klein     & Klein     & Groot        \\
		Focus-effect                 & Geen      & Scherm dimt & Geen        \\ \midrule
	\end{tabular}
	\begin{flushleft}
		\footnotesize *Ondersteunt extra opties zoals tekstkleur, hyperlinks en lijst-opmaak.
	\end{flushleft}
\end{table}

\subsubsection*{2. Walkthroughs}
Waar een prompt beperkt is tot één enkele interactie op één scherm, biedt een walkthrough de functionaliteit om een complexer proces te begeleiden, waar meerdere stappen na elkaar gekoppeld zijn.

\begin{itemize}
	\item \textbf{Prompts:} Een walkthrough is in essentie een aaneenschakeling van maximaal tien afzonderlijke prompts (targeted, floating of docked) die een logische volgorde vormen.
	\item \textbf{Navigatie:} Walkthroughs zijn in staat om gebruikers over verschillende webpagina's, tabbladen en zelfs verschillende sub-applicaties binnen Salesforce heen te leiden.
	\item \textbf{Contextueel:} De kracht van een walkthrough ligt in de timing. De informatie is direct gekoppeld aan de handeling van de gebruiker op het scherm. Dit voorkomt dat een gebruiker vastloopt of fouten maakt, omdat de juiste instructies precies op het moment en de plaats komen waar deze gebruiker er nood aan heeft.
\end{itemize}

\subsubsection*{3. Configuratie en Gedrag}
De effectiviteit van in-app guidance wordt onder meer bepaald door de instellingen hiervan. De prompts en walkthroughs kunnen geconfigureerd worden in een guidance builder waar onder andere de inhoud, locatie en type gekozen kunnen worden.

\begin{itemize}
	\item \textbf{Locatie en Context:} Per prompt wordt de reikwijdte bepaald. Dit varieert van een specifieke pagina binnen één applicatie tot een globale weergave die overal in Salesforce verschijnt (Any Page, Any App). Dit maakt het mogelijk om de guidance precies af te stemmen op de specifieke omgeving waarin de gebruiker op dat moment werkt.
	
	\item \textbf{Permissies:} De zichtbaarheid van de in-app guidance kan worden gefilterd op basis van profielen en permissies. Hierdoor is het mogelijk om in-app guidance enkel te tonen aan gebruikers die over de juiste rechten beschikken om de betreffende acties uit te voeren.
	
	\item \textbf{Activatie en Herhaling:} Voor zowel losse prompts als volledige walkthroughs kan worden bepaald wanneer ze voor het eerst verschijnen. Dit kan op twee manieren gebeuren, namelijk op basis van een vooraf ingesteld schema, of door een actie van de gebruiker. Het kan bijvoorbeeld dagelijks verschijnen op een specifieke pagina, of de gebruiker klikt op een knop die de guidance start. Hierbij kan de frequentie van herhaling vooraf ingesteld worden, en de gebruiker kan deze ook altijd aanpassen.
	
	\item \textbf{Interactie en Media:} Een prompt kan worden uitgebreid met actieknoppen die verwijzen naar een (interne) URL of flow. Daarnaast kan de tekst visueel worden ondersteund met afbeeldingen of, in het geval van docked prompts, met video's die direct in de interface kunnen worden afgespeeld.
\end{itemize}

\subsection{Evaluatie}
Naast de creatie en configuratie van in-app guidance biedt Salesforce enkele middelen om het gebruik ervan te analyseren:

\begin{itemize}
	\item \textbf{Gebruikersinzichten (Analytics):} Het systeem verzamelt automatisch data over de interactie met de guidance, zoals het aantal weergaven, het gebruik van actieknoppen (\textit{Click-through Rate}) en het aantal keer dat een prompt of stap voortijdig wordt gesloten. Deze inzichten stellen admins in staat om te beoordelen of de getoonde informatie relevant is en waar instructies eventueel geoptimaliseerd moeten worden.
	
\end{itemize}

\subsection{Restricties en Limitaties}
Zowel Salesforce zelf als de native In-App Guidance functionaliteit kent enkele strikte platformlimieten waarmee tijdens de implementatie rekening gehouden moet worden:

\begin{itemize}
	\item \textbf{Desktop Exclusief:} In-App Guidance is exclusief ontworpen voor de desktopversie van Lightning Experience; de functionaliteit wordt momenteel niet ondersteund in de mobiele Salesforce-applicatie.
	\item \textbf{Systeeminstellingen en Custom Schermen:} Prompts kunnen niet worden geplaatst op systeeminstellingen (Setup-pagina's) of op specifieke onderdelen van de interface die buiten het standaard Lightning-framework vallen.
	\item \textbf{Stappenlimiet:} Een enkele walkthrough is technisch gelimiteerd tot maximaal tien opeenvolgende stappen (prompts).
	\item \textbf{Governor Limits:} Omdat Salesforce een \textit{multi-tenant} platform is (waarbij servercapaciteit wordt gedeeld tussen klanten), gelden er harde limieten voor server-side processen. Apex-logica is strikt beperkt in het aantal SOQL-queries en DML-operaties per transactie.
	\item \textbf{UI-Gedrag (Lightning Web Security):} Salesforce hanteert een strikte beveiligingsarchitectuur (\textit{Lightning Web Security} ) die voorkomt dat componenten buiten hun eigen grenzen treden. Hierdoor is het voor een custom component technisch onmogelijk om standaard Salesforce UI-gedrag programmatisch te overschrijven.
\end{itemize}

\section{Requirements Analyse}%
\label{sec:Requirements analyse}

Uit de voorgaande Salesforce analyse lijkt de native In-App Guidance uitstekend geschikt om een gebruiker te begeleiden bij allerlei soorten processen. Voor de initiële setup van een complexe applicatie zoals 24flow moet deze guidance ook duidelijk georganiseerd worden. Daarom wordt er voor het proof of concept gewerkt met een op maat gemaakt 'Onboarding Dashboard' waar de gebruiker op landt bij het starten van de applicatie. Dit dashboard fungeert als een centrale hub die de gebruiker een duidelijk stappenplan en houvast biedt. Vanuit dit overzicht kan de gebruiker dan stap per stap de nodige configuraties opzetten met behulp van de in-app guidance zonder het overzicht te verliezen en met een duidelijke houvast.
\\

Voor de technische uitwerking van deze oplossing is het cruciaal om de vereisten van het systeem helder in kaart te brengen. In deze sectie worden de kerncomponenten van 24flow, de doelgroep, de niet-functionele requirements en de specifieke functionele requirements gedefinieerd en geprioritiseerd.
\subsection{Kerncomponenten van 24flow}
\label{subsec:kerncomponenten_24flow}

Om de complexiteit van de onboarding te begrijpen, is inzicht in de modulaire opbouw van 24flow noodzakelijk. Voor dit onderzoek wordt er enkel gefocust op één van de verschillende modules die 24flow aanbiedt, namelijk 24factory, waar er een digitale fabriek gebruikt wordt die de effectieve werkvloer beheert van een klant. Dit platform bestaat uit verschillende interactieve componenten die nauw met elkaar verbonden zijn via Salesforce-data:

\begin{itemize}
	\item \textbf{Order Pool:} Een centraal overzicht waar alle productieorders binnenkomen. Deze component fungeert als de ``ingang'' van de fabriek. De effectiviteit van de Order Pool hangt af van de juiste filtering en weergave van velden (zoals leverdatum en producttype). 
	\item \textbf{Team Cockpit:} Een visueel bord waarop operatoren de status van Productie Orders (bijv. \textit{In Buffer, In Process, Expected}) kunnen opvolgen. Dit is het hart van de werkvloer planning, en hiermee kan de hoeveelheid work in progress op de vloer georganiseerd worden.
	\item \textbf{Production Orders, Teams \& Operations:} Dit zijn enkele van de belangrijkste achterliggende datastructuren. Een productieorder bestaat uit meerdere bewerkingen (operations) die door verschillende teams uitgevoerd moeten worden.
\end{itemize}

\subsubsection*{De afhankelijkheid van data}
Een fundamenteel kenmerk van deze componenten is dat zij \textbf{visueel georiënteerd} zijn. Een Order Pool zonder records is een lege tabel, en een Team Cockpit zonder actieve bewerkingen toont enkel lege kolommen.
\\

Daarom is er de noodzaak om een gebruiker over realistische data te laten beschikken tijdens de onboarding. Ofwel zou er demo data gebruikt kunnen worden of een klant zou zijn eigen data kunnen uploaden en gebruiken of een combinatie van beide indien mogelijk. Dit zou ervoor zorgen dat de gebruiker direct met een gevulde interface kan werken, wat essentieel is voor een succesvolle configuratie.

\subsection{Doelgroep}
\label{subsec:doelgroepanalyse}

Voor de uitwerking van een effectieve onboarding-oplossing is het cruciaal om te bepalen voor wie het systeem wordt ontworpen. Binnen het 24flow-ecosysteem zijn er twee specifieke gebruikersrollen geïdentificeerd. Hoewel de uiteindelijke applicatie voor beide rollen hetzelfde is, verschillen hun doelen en vereisten:

\begin{itemize}
	\item \textbf{De Salesforce-beheerder van de klant:} Dit is een beheerder bij een (potentieel) nieuw klantbedrijf die 24flow wil evalueren of implementeren. Deze gebruiker kent de applicatie nog niet en wil het systeem zelfstandig leren gebruiken en testen. Voor deze groep is er een sterke behoefte aan stapsgewijze, interactieve begeleiding en direct toegankelijke en duidelijke documentatie. Daarnaast is het voor hen belangrijk om over realistische testdata te beschikken, zodat zij de applicatie direct visueel kunnen evalueren zonder eerst uren aan handmatige data-invoer te spenderen. Het uiteindelijke doel van deze gebruiker is 24flow in werking zien en indien ze nog geen klant zijn evalueren of ze wel of niet klant willen worden.
	
	\item \textbf{De Solution Engineer (24flow-expert):} Dit is een interne medewerker van 24flow. Hoewel het onboarding-systeem primair bedoeld is voor (nieuwe) klanten, bieden specifieke onderdelen ervan ook een grote meerwaarde voor deze interne experts. Zij kennen de applicatie al door en door en hebben geen behoefte aan stapsgewijze handleidingen. Echter, de functionaliteiten die nodig zijn om een omgeving geautomatiseerd op te zetten en van testdata te voorzien, kunnen door hen uitstekend worden ingezet om in zeer korte tijd bijvoorbeeld een demo-omgeving klaar te zetten.
\end{itemize}

\subsection{Centralisatie}
\label{subsec:centralisatie_dashboard}

Een van de grootste uitdagingen bij een standaard Salesforce-implementatie is de verspreiding van de instellingen. Omdat de benodigde configuraties verspreid liggen over talloze menu's, objecten en Setup-pagina's, ontstaat er al snel een gefragmenteerde gebruikerservaring. Om te voorkomen dat een nieuwe gebruiker het overzicht verliest, is er een duidelijke noodzaak voor centralisatie. 
\\

Het te ontwikkelen systeem moet daarom een centraal vertrekpunt bevatten dat voor structuur bij de eindgebruiker kan zorgen. In plaats van de gebruiker te verplichten zelf op zoek te gaan naar alle componenten, moeten alle noodzakelijke onboarding-taken op deze centrale plaats worden samengebracht. Dit vermijdt onnodig zoekwerk en zorgt ervoor dat een gebruiker het overzicht steeds bewaart. Hier kan dan ook een link naar de documentatie zijn en dit kan eveneens alle extra tools bevatten die nodig zouden kunnen zijn voor een goede onboarding.

\subsection{Niet-Functionele Requirements}
De oplossing moet voldoen aan de volgende technische randvoorwaarden binnen het Salesforce-platform:

\begin{itemize}
	\item \textbf{Interface:} De gebruikersinterface moet gebouwd worden conform het Salesforce Lightning Design System (SLDS).
	\item \textbf{Performantie:} Alle backend-processen (zoals Apex data-generatie) moeten binnen de vastgestelde \textit{governor limits} van Salesforce opereren.
	\item \textbf{Beheer:} De configuratie van de in-app guidance (prompts en walkthroughs) moet aanpasbaar zijn via de standaard Salesforce set-up, zonder dat er codewijzigingen nodig zijn.
\end{itemize}

\subsection{Prioritering van de requirements}
Om de scope van het proof of concept onder controle te houden en te garanderen dat de meest kritische onderdelen als eerste worden gebouwd, zijn de verzamelde requirements geprioriteerd volgens de MoSCoW-methode.

\subsubsection*{Must haves}
De "Must haves" vormen de minimale set aan functionaliteiten die aanwezig moeten zijn om een werkbare oplossing te bieden die beantwoordt aan de onderzoeksvraag:

\begin{itemize}
	\item \textbf{Zelfstandige Systeemconfiguratie:} Een volledige interactieve onboarding die de gebruiker stap-voor-stap door de initiële opzet van 24flow loodst voor zijn belangrijkste componenten. Dit stelt de gebruiker in staat om de applicatie volledig operationeel te krijgen zonder externe hulp.
	
	\item \textbf{Centraal Dashboard:} Een overzichtelijk startpunt waar de gebruiker de voortgang van zijn configuratie kan monitoren via een visuele voortgangsbalk en status-indicatoren. Dit is eveneens de landingspagina van elke 24flow applicatie.

	\item \textbf{Directe Navigatie en Activatie:} Een koppeling waarbij de gebruiker met één druk op de knop direct naar de juiste pagina wordt gestuurd en de bijbehorende  walkthrough automatisch wordt gestart.
	
	\item \textbf{Automatische Voortgangsdetectie:} Het systeem moet zelfstandig herkennen wanneer een taak succesvol is afgerond en dit direct verwerken in het dashboard. Het moet per onderdeel bijhouden of en in hoeverre de configuratie is voltooid.
	
	\item \textbf{Demo Data Generator:} De mogelijkheid om de omgeving te vullen met actuele testdata. Dit is essentieel voor een realistisch gebruik van de applicatie, zodat de gebruiker de onboarding doorloopt in een "gevulde" interface.
	
	\item \textbf{Hiërarchische Taakstructuur:} De mogelijkheid om grote, complexe processen op te delen in kleinere subtaken, zodat de gebruiker stap-voor-stap door de configuratie wordt geleid zonder het overzicht te verliezen.
	
	\item \textbf{Navigatie tussen objecten:} De mogelijkheid voor walkthroughs om de gebruiker naadloos over verschillende tabbladen en records heen te leiden, aangezien deze zelden beperkt blijven tot één enkel scherm.
	
	\item \textbf{Documentatie:} Een directe en permanente link naar de officiële documentatie op het dashboard. Dit dient als hulpmiddel en naslagwerk voor gebruikers die specifieke, of meer diepgaande informatie zoeken over bepaalde componenten.


\end{itemize}
\subsubsection*{Should haves}
	De "Should haves" zijn belangrijke functionele aanvullingen die de kwaliteit van de onboarding verhogen en essentieel zijn voor een volledige en professionele gebruikerservaring:

\begin{itemize}
	\item \textbf{Video-integratie:} Het toevoegen van korte video-tutorials in docked prompts voor de meest complexe processen. Dit verhoogt de effectiviteit van de onboarding, maar kan in eerste instantie worden vervangen door tekst en afbeeldingen.
	
	\item \textbf{Data Reset-functie:} Een veilige methode om testdata te wissen, zodat een gebruiker altijd met een propere omgeving kan werken.
	
	\item \textbf{Opladen Eigen Data:} Een methode om eigen data te uploaden die gebruikt kan worden voor de onboarding. Dit kan in eerste instantie een placeholder zijn die later kan worden geïmplementeerd.
	
	\item \textbf{Taakbevestiging:} Voor stappen die niet automatisch door het systeem gemeten kunnen worden, moet de gebruiker zelf een taak als voltooid kunnen markeren.

	\item \textbf{Custom Branding:} Het aanpassen van de kleuren van de prompts aan de huisstijl van 24flow. Dit draagt bij aan de uniformiteit en herkenbaarheid binnen het platform.
	
	\item \textbf{Documentatie Linken} Getoonde documentatie moet specifiek gekoppeld zijn aan de respectievelijke component. Een gebruiker mag niet op zoek moeten gaan in alle documentatie voor een specifiek onderdeel.
	
\end{itemize}

\subsubsection*{Could haves}
Deze elementen worden als "extra" beschouwd en worden enkel geïmplementeerd als er na de kritieke onderdelen nog tijd over is:

\begin{itemize}
	\item \textbf{Hervatten:} De optie voor gebruikers om een reeds voltooide walkthrough op elk moment opnieuw te starten en verder te gaan op het punt waar ze gebleven zijn.
	
	\item \textbf{Gebruikersstatistieken (Analytics):} Het opzetten van rapportage om te meten hoe vaak walkthroughs worden voltooid en waar gebruikers vastlopen. Dit is nodig voor de evaluatie op lange termijn, maar niet voor de technische werking.
	
	\item \textbf{Gebruikersfeedback:} De mogelijkheid voor gebruikers om direct binnen een prompt aan te geven of de instructie nuttig was via een simpele duimpje-omhoog/omlaag interactie.
	
	\item \textbf{Gamificatie-elementen:} Het toevoegen van visuele elementen, zoals een vuurwerk animatie of een digitale "badge" bij het voltooien van een bepaalde configuratie, om de gebruikerservaring te verbeteren
	
	\item \textbf{Zoekfunctionaliteit:} Een zoekbalk op het dashboard waarmee gebruikers direct naar specifieke instellingen of documentatie-onderwerpen kunnen springen.
	
	\item \textbf{Directe Support-link:} Een knop op het dashboard die direct een support-ticket opent bij een 24flow-expert voor vragen die niet door de onboarding worden beantwoord.
\end{itemize}

\subsubsection*{Won't haves}
Deze onderdelen zijn expliciet buiten de scope van dit project gelaten, vaak vanwege technische beperkingen van het Salesforce-platform:

\begin{itemize}
	\item \textbf{Mobiele Ondersteuning:} Ondersteuning voor de Salesforce Mobile App, aangezien in-app guidance momenteel technisch enkel beschikbaar is voor de desktopomgeving.

	\item \textbf{Chatbot:} AI Chatbots: Een interactieve bot die vragen beantwoordt; we focussen op gestuurde navigatie, de gebruiker moet in staat zijn om zelf wegwijs te worden op het platform.
\end{itemize}

\section{Ontwerp proof of concept}
\label{sec:Ontwerp}
\subsection{Use cases}

\subsubsection{Use Case: Initiële configuratie van de Order Pool}

Deze use case beschrijft hoe een nieuwe gebruiker zelfstandig een essentieel onderdeel van de applicatie inricht via de onboarding oplossing.
\\

Actor: Nieuwe 24flow-gebruiker (Systeembeheerder/Admin).
\\

Doel: Het succesvol en zelfstandig configureren van de Order Pool-instellingen.

\begin{enumerate}
	\item \textbf{Startpunt:} De gebruiker landt op de centrale landingspagina van de applicatie. Het \textbf{Onboarding Dashboard} verschijnt direct in beeld.
	
	\item \textbf{Status:} De gebruiker ziet in de voortgangsbalk dat de setup op 0 procent staat. In de hiërarchische takenlijst wordt de taak "Order Pool instellen" als openstaand gemarkeerd.
	
	\item \textbf{Activatie:} De gebruiker klikt op de actieknop naast de taak. Het systeem navigeert de gebruiker direct naar de juiste instellingenpagina van de Order Pool.

	\item \textbf{Begeleiding:} Bij het laden van de pagina start automatisch de Salesforce \textbf{Walkthrough}. Een targeted prompt wijst naar het eerste configuratieveld en legt uit welke informatie hier benodigd is.

	\item \textbf{Interactie:} De gebruiker volgt de stappen van de walkthrough, waarbij hij door verschillende tabbladen wordt geleid om filters en kolommen te definiëren.

	\item \textbf{Detectie:} Nadat de gebruiker op 'Opslaan' heeft geklikt, herkent de \textbf{Automatische Voortgangsdetectie} op de achtergrond dat de configuratie-records zijn aangemaakt.

	\item \textbf{Voltooiing:} De gebruiker keert terug naar het dashboard. De taak "Order Pool instellen" is nu voorzien van een groen vinkje en de algehele voortgangsbalk is bijgewerkt naar een hoger percentage.
\end{enumerate}

\subsubsection{Use Case: Beheer van testdata via het dashboard}
Deze use case illustreert hoe de functionaliteiten in de zijbalk van het dashboard worden ingezet om de applicatie klaar te maken voor gebruik door testdata te genereren.
\\

Actor: Nieuwe 24flow-gebruiker (Systeembeheerder/Admin).
\\

Doel: Het succesvol vullen van de applicatie met realistische testdata, zodat de functionaliteiten optimaal kunnen worden geconfigureerd en gebruikt.

\begin{enumerate}
	\item \textbf{Startpunt:} De gebruiker bevindt zich op het Onboarding Dashboard en merkt op dat de applicatie nog leeg is en geen data bevat.
	\item \textbf{Actie:} In het "Tools" paneel aan de rechterkant klikt de gebruiker op de knop "Generate Demo Data".
	\item \textbf{Verwerking:} Het systeem toont tijdelijk een laad-indicator. Op de achtergrond wordt een Apex-proces aangeroepen dat automatisch een set gestandaardiseerde testrecords aanmaakt.
	\item \textbf{Bevestiging:} Zodra het proces is afgerond, toont het dashboard een succesmelding (toast notificatie) aan de gebruiker.
	\item \textbf{Resultaat:} De applicatie is nu gevuld met data, waardoor de gebruiker bij walkthroughs direct deze data kan gebruiken om de nodige configuraties te doen.

\end{enumerate}



\clearpage
\subsection{Wireframes}

Om de requirements visueel te vertalen, zijn er enkele wireframes uitgewerkt. Dit ontwerp dient als blauwdruk voor de uiteindelijke Lightning Web Component en de onboarding implementatie.


\subsubsection{Wireframe: Onboarding Dashboard}


Zoals te zien is in Figuur \ref{fig:wireframe_dashboard}, fungeert het dashboard als de landingspagina en het vertrekpunt voor een nieuwe gebruiker. Aan de linkerkant bevindt zich de hiërarchische taakstructuur met dynamische status-indicatoren en een voortgangsbalk. Aan de rechterkant zijn enkele tools om bijvoorbeeld data te genereren en wissen en documentatie te raadplegen.
\begin{figure}[htbp]
	\centering
	% Gebruik het relatieve pad en zorg dat de backslashes ( \ ) vervangen zijn door forward slashes ( / )
	\includegraphics[width=0.95\textwidth]{WireframeDashboardOnboarding}
	\caption{Wireframe van het Onboarding Dashboard inclusief voortgangsbalk en zijbalk met extra tools.}
	\label{fig:wireframe_dashboard}
\end{figure}



\clearpage
\subsubsection{Wireframe: Walkthrough Stap}
Wanneer de gebruiker in het dashboard een taak start, navigeert het systeem direct naar de desbetreffende configuratiepagina. Figuur \ref{fig:wireframe_walkthrough} toont deze stap. Hier neemt de in-app guidance het proces over in de vorm van een walkthrough. Een \textit{Targeted Prompt} wijst de gebruiker expliciet op het in te vullen veld en biedt direct een aanbeveling ("Action Link") voor meer info.
\\

 Tegelijkertijd biedt een \textit{Docked Prompt}, onderaan het scherm, extra ondersteuning via een instructievideo en een documentatie link. Dit ontwerp zorgt ervoor dat de gebruiker complexe acties stapsgewijs en zelfstandig kan uitvoeren en ook steeds een directe link heeft naar extra documentatie indien nodig.

\begin{figure}[htbp]
	\centering
	% Gebruik het relatieve pad en zorg dat de backslashes ( \ ) vervangen zijn door forward slashes ( / )
	\includegraphics[width=0.95\textwidth]{WalkthroughStep}
	\caption{Wireframe van een walkthrough stap voor het configureren van de Order Pool.}
	\label{fig:wireframe_walkthrough}
\end{figure}

\section{Realisatie en Implementatie}
\label{sec:realisatie}
In dit hoofdstuk wordt de technische en functionele realisatie van het proof of concept besproken. Waar de voorgaande secties zich richtten op de vereisten en het conceptuele ontwerp, focust dit deel op de implementatie aan de hand van de vooropgestelde requirements en wireframes. 

De architectuur van de oplossing is fundamenteel opgebouwd uit twee samenwerkende lagen. Enerzijds is er de Salesforce In-App Guidance, die verantwoordelijk is voor de actieve begeleiding van de gebruiker over verschillende schermen en stappen heen. Anderzijds is er een setup assistant applicatie ontwikkeld: het Onboarding Dashboard. Dit dashboard, gebouwd als een moderne Lightning Web Component (LWC), fungeert als het controlecentrum. Het verdeelt de onboarding over kleinere stappen, stuurt de gebruiker naar de juiste onderdelen, en valideert op de achtergrond (via Apex) of de nodige data en configuraties succesvol zijn aangemaakt.

In de volgende paragrafen worden de belangrijkste pijlers van deze implementatie toegelicht: de dynamische opbouw van het dashboard, de integratie met de Salesforce guidance tools en de logica achter de voortgangsdetectie.

\subsection{Architectuur van de front-end}

De visuele kern van het proof of concept wordt dus gevormd door het 'Onboarding Dashboard'. Omdat de nodige configuraties verspreid zijn over allerlei schermen, menu's en salesforce instellingen, is er gekozen om dit dashboard als centrum voor de onboarding te laten dienen. De gebruiker vertrekt steeds vanuit deze pagina om alles op te zetten en wordt ook steeds hiernaartoe teruggeleid.

\subsubsection*{Opbouw van de User Interface (Dashboard)}
De interface is volledig gebouwd met Lightning Web Components (LWC) en het Salesforce Lightning Design System (SLDS). Het benutten van deze standaard architectuur is om twee redenen cruciaal voor de effectiviteit van het dashboard:

\begin{enumerate}
	\item \textbf{Consistente User Experience:} Door SLDS toe te passen, is de component visueel niet te onderscheiden van de standaard Salesforce-omgeving. Dit vergroot het vertrouwen en de herkenbaarheid voor de gebruiker.
	\item \textbf{Reactieve UI:} LWC leunt sterk op \textit{reactive data binding}. Wanneer de backend (Apex) registreert dat een configuratiestap is voltooid, of wanneer een proces loopt dat even moet laden (zichtbaar via de \texttt{lightning-spinner} component), worden de interface-elementen direct en asynchroon geüpdatet, zonder dat de gebruiker de pagina hoeft te herladen.
\end{enumerate}

Initieel in de wireframes was het dashboard zo opgebouwd dat aan de linkerkant de configuratiestappen stonden en rechts enkele extra tools.
Om de gebruiker echter zoveel mogelijk op het centrale scherm te houden en de documentatie hier ook rechtstreeks te kunnen tonen, is het dashboard in de uiteindelijke implementatie visueel opgedeeld in twee secties: een linkerkolom met de echte configuratiestappen en enkele tools daarboven en een rechterkolom puur om documentatie te tonen.
Initieel in de wireframes was het dashboard zo opgebouwd dat aan de linkerkant de configuratiestappen stonden en rechts enkele extra tools. Om de gebruiker echter zoveel mogelijk op het centrale scherm te houden en de documentatie hier ook rechtstreeks te kunnen tonen, is het dashboard in de uiteindelijke implementatie visueel opgedeeld in twee secties: een linkerkolom met de echte configuratiestappen en enkele tools daarboven en een rechterkolom puur om documentatie te tonen.

\begin{figure}[hbt!]
	\centering
	\includegraphics[width=0.9\textwidth]{Dashboard.png}
	\caption{Overzicht van het Onboarding Dashboard}
	\label{fig:dashboard}
\end{figure}

\vspace{0.3cm}
\noindent\textbf{A. De Setup-kolom (Configuratie \& Tools)}\\
Aan de linkerkant bevindt zich de interactieve onboarding component (\textit{24flow Setup Assistant}). Deze is opgebouwd uit drie logische lagen:
\begin{itemize}
	\item \textbf{Voortgang, Data \&  Default configuratie:} Bovenaan bevindt zich een overkoepelende voortgangsbalk, waarop de gebruiker in percentages ziet hoeveel hij al geconfigureerd heeft. Deze balk is reactief: de logica combineert actuele database-controles uit de backend met lokale browser-opslag (\textit{localStorage}) voor handmatige taken, om zo het percentage te berekenen. Direct daaronder worden twee blokken getoond met algemene configuratie: 
	\begin{itemize}
		\item \textit{Demo Data Management:} Enerzijds een knop om via een Apex-proces automatisch demodata te genereren, zodat de gebruikersomgeving gevuld wordt met realistische productiedata. Anderzijds een placeholder knop om een eigen CSV-bestand in te laden met echte data van de klant, dit is een placeholder omdat het omvormen van deze data naar de Salesforce-omgeving buiten de scope van deze proof of concept valt. Tijdens deze zware database-acties vergrendelt de component de interface met een laadscherm (\textit{spinner}) om dubbele acties te voorkomen. Onderaan de component is eveneens een reset-knop voorzien om deze data veilig te verwijderen en de gebruikersomgeving op te kuisen.
		\item \textit{Master Auto-Configuration:} Een geavanceerde functionaliteit waarmee de gebruiker voor alle 24flow componenten een aantal default instellingen aanmaakt waar mogelijk. Dit is bedoeld voor de gebruiker die snel 24flow in werking wil zien en wil kunnen testen zonder hiervoor alle configuratie in detail te moeten doorlopen.
	\end{itemize}
	\item \textbf{Hiërarchische Takenlijst:} Daaronder volgt het stappenplan in een uitklapbare (accordion) structuur. Configuratietaken zijn gegroepeerd met duidelijke titels, beschrijvingen, een indicatie van de geschatte uitvoeringstijd en een aanduiding of deze stap wel of niet gedaan is. Hierop bevindt zich ook een knop om voor deze stap automatisch een default configuratie te laten maken.
	\item \textbf{Substappen:} Binnen elke hoofdtaak bevinden zich enkele specifieke substappen met een geschatte uitvoeringstijd. Op dit niveau kan de gebruiker ervoor kiezen om een stap handmatig te configureren (wat de navigatie naar de juiste pagina triggert, waarbij het systeem op de achtergrond een signaal meegeeft om direct de juiste In-App Guidance walkthrough op te starten), direct gerelateerde documentatie te openen, of de taak handmatig als voltooid te markeren. Dit laatste is nodig omdat niet automatisch gecontroleerd kan worden of bepaalde salesforce setup instellingen en salesforce pagina's zijn aangemaakt.
\end{itemize}
\begin{figure}[hbt!]
	\centering
	\includegraphics[width=0.9\textwidth]{DashboardOpen.png}
	\caption{Uitgeklapte hiërarchische takenlijst binnen het dashboard.}
	\label{fig:dashboard_open}
\end{figure}

\vspace{0.3cm}
\noindent\textbf{B. Documentatie-viewer}\\
Aan de rechterkant van het scherm bevindt zich de documentatie-viewer. Technisch gezien is deze component opgebouwd rond een \texttt{iframe}, waarmee de externe webpagina's van de 24flow-documentatie direct binnen de Salesforce-interface worden getoond. Door deze direct naast de setup-kolom te plaatsen, kan de gebruiker de uitleg uit de 24flow documentatie raadplegen en tegelijkertijd links in het scherm de configuratie-knoppen bedienen.
\\ 

Technisch communiceren de twee kolommen via de \textit{Lightning Message Service} (LMS); zodra er links een help-actie plaatsvindt, ontvangt de viewer een bericht en laadt direct de juiste URL in het \texttt{iframe}. Hiermee kan de gebruiker per configuratie stap direct de gelinkte documentatie openen in hetzelfde scherm. Dit elimineert de noodzaak om constant tussen verschillende browsertabbladen te wisselen en zelf op zoek te gaan naar de locatie van bepaalde documentatie, wat het vinden van de nodige info aanzienlijk versnelt.
\\

Om de documentatie mee te kunnen nemen doorheen de rest van de applicatie en niet enkel op het initiële scherm te tonen is er in de Salesforce utility bar aan de onderkant van het scherm steeds een knop met een venster waarin de Documentatie-Viewer ook te zien is. Hierbij is de effectief getoonde documentatie rechtstreeks gelinkt met de pagina waarop de gebruiker zich bevindt, deze is dus context-specifiek.
\begin{figure}[hbt!]
	\centering
	\includegraphics[width=0.9\textwidth]{DocumentationViewer.png}
	\caption{De documentatie-viewer geopend via de Salesforce Utility Bar.}
	\label{fig:doc_viewer}
\end{figure}

\vspace{0.3cm}
\noindent\textbf{C. Interactieve Begeleiding (Walkthroughs)}\\
De visuele ondersteuning wordt afgemaakt door de integratie van Salesforce In-App Guidance. Waar het dashboard en de viewer het controlecentrum vormen, vindt de effectieve begeleiding plaats op de verschillende applicatiepagina's zelf.
\\

Zodra de gebruiker een taak start vanuit het dashboard, wordt de gebruiker verplaatst naar de juiste pagina voor die bepaalde configuratiestap. Hierbij start de gekoppelde walkthrough en is de eerste prompt zichtbaar, dit is ofwel een targeted prompt of een floating prompt afhankelijk van de situatie. Aan de hand van deze gekoppelde prompts wordt de gebruiker dan stap voor stap gestuurd bij het opzetten van zijn instellingen. De laatste prompt bevat steeds een link terug naar het dashboard en eventueel is er ook een docked prompt aanwezig met een video of andere multimedia indien dit beschikbaar is.

\begin{figure}[hbt!]
	\centering
	\includegraphics[width=0.9\textwidth]{WalktroughStepImpl.png}
	\caption{Start van een walkthrough met een targeted prompt (stap 1).}
	\label{fig:walkthrough_1}
\end{figure}

\begin{figure}[hbt!]
	\centering
	\includegraphics[width=0.9\textwidth]{WalktroughStepImpl2.png}
	\caption{Vervolgstap binnen de walkthrough met een floating prompt (stap 2).}
	\label{fig:walkthrough_2}
\end{figure}

\subsection{Architectuur van de Backend}

De backend van het proof of concept heeft twee duidelijke kerntaken: het aanleveren van de inhoud voor het dashboard en het uitvoeren van de verschillende dashboard acties in de achtergrond.

\subsubsection*{Datastructuur: De Opbouw van het Dashboard}
\label{subsec:dynamische_datastructuur}

Om het dashboard overzichtelijk te houden, zitten de teksten, stappen en documentatielinks niet vast in de front-end code. In plaats daarvan wordt de structuur van het onboardingproces bewaard als data in een \textit{Custom Metadata Type}.
\\

Bij het inladen roept het dashboard deze configuratie op. De applicatie leest deze data (geformatteerd als JSON) uit om de stappenlijst op het scherm visueel op te bouwen. Door de inhoud en de weergave op deze manier te scheiden, is het mogelijk om later eenvoudig stappen toe te voegen of aan te passen zonder dat de code van de interface gewijzigd hoeft te worden.

\subsubsection*{Uitvoerende Logica: De Backend Controller}
\label{subsec:backend_controller}

Terwijl de front-end zorgt voor de visuele weergave, regelt een centrale Apex-klasse (\texttt{OnboardingController}) de effectieve bedrijfslogica en de communicatie met de database. Deze backend heeft vijf concrete hoofdtaken:

\begin{itemize}
	\item \textbf{Data Ophalen:} Het inladen van de dashboard-opbouw vanuit de metadata, zodat de stappenlijst direct correct op het scherm verschijnt.
	\item \textbf{Statuscontrole:} De backend controleert met gerichte database-queries of taken daadwerkelijk zijn uitgevoerd. Hierdoor toont de voortgangsbalk altijd de actuele, ware status, ongeacht wat er in de interface wordt aangeklikt.
	\item \textbf{Automatische Configuratie:} Configuraties kunnen met één actie worden geautomatiseerd. Het systeem genereert automatisch enkele default waarden voor de verschillende instellingen.
	\item \textbf{Demodata Beheren:} Om te voorkomen dat gebruikers in een lege applicatie starten, kan de backend realistische testdata genereren. Daarnaast bevat de controller de logica om de gebruikersomgeving weer op te kuisen door alle gegenereerde testdata en instellingen in één keer veilig te verwijderen.
	\item \textbf{Acties Loggen:} Belangrijke handelingen op het dashboard worden opgeslagen in een log-object. Dit maakt het mogelijk om achteraf te analyseren hoe gebruikers het onboardingproces doorlopen.
\end{itemize}
\section{Evaluatie}%
\label{sec:Evaluatie}

In dit hoofdstuk wordt het ontwikkelde proof of concept kritisch geëvalueerd. Het doel van deze evaluatie is om vast te stellen in hoeverre de hybride onboarding-oplossing beantwoordt aan de vooraf gestelde requirements, hoe het systeem presteert onder praktische scenario's, en waar zich de grenzen van de huidige implementatie bevinden.

\subsection{Toetsing aan de Requirements (MoSCoW)}
De rode draad van deze evaluatie is de toetsing aan de requirements-analyse. Onderstaande tabel geeft een volledig overzicht van alle geprioriteerde vereisten en hun status in het uiteindelijke PoC.

\begin{table}[h!]
	\centering
	\caption{Volledig statusoverzicht van de MoSCoW-requirements}
	\label{tab:moscow_volledig}
	\small
	\begin{tabular}{@{}llc@{}}
		\toprule
		\textbf{Prioriteit} & \textbf{Requirement} & \textbf{Status} \\ \midrule
		\textbf{Must Have} & Zelfstandige Systeemconfiguratie (Onboarding flow) & Voltooid \\
		& Centraal Dashboard (LWC Hub) & Voltooid \\
		& Directe Navigatie en Walkthrough Activatie & Voltooid \\
		& Automatische Voortgangsdetectie (Apex validation) & Voltooid \\
		& Demo Data Generator (Apex) & Voltooid \\
		& Hiërarchische Taakstructuur (Geneste JSON) & Voltooid \\
		& Navigatie tussen objecten (Cross-tab/App) & Voltooid \\
		& Documentatie Link op Dashboard & Voltooid \\ \midrule
		\textbf{Should Have} & Video-integratie (Docked Prompts) & Voltooid \\
		& Data Reset-functie (Clear all data) & Voltooid \\
		& Opladen Eigen Data (CSV Upload) & Placeholder \\
		& Taakbevestiging (Manueel markeren via LocalStorage) & Voltooid \\
		& Custom Branding (SLDS \& Huisstijl) & Voltooid \\
		& Documentatie Contextueel Linken (LMS per component) & Voltooid \\ \midrule
		\textbf{Could Have} & Hervatten (Restart-logica \& Statusbehoud) & Voltooid \\
		& Gebruikersstatistieken (Telemetrie via \texttt{logAction}) & Voltooid \\
		& Gebruikersfeedback (Duimpjes) & Niet uitgevoerd \\
		& Gamificatie (Animaties/Badges) & Niet uitgevoerd \\
		& Zoekfunctionaliteit / Directe Support-link & Niet uitgevoerd \\ \midrule
		\textbf{Won't Have} & Mobiele Ondersteuning \& Chatbot & Buiten scope \\ \bottomrule
	\end{tabular}
\end{table}


\subsubsection*{Analyse van de 'Must Haves'}
Alle kritieke vereisten zijn succesvol geïmplementeerd en vormen de fundering van het systeem. De belangrijkste realisaties zijn:

\begin{itemize}
	\item \textbf{Centraal Dashboard en Slimme Navigatie:} De LWC-hub fungeert als de centrale 'controlekamer' van waaruit alle acties vertrekken. De navigatie-logica gaat verder dan simpele links door het injecteren van een \textit{promptId} in de URL-parameters. Hierdoor start de juiste Salesforce Walkthrough automatisch op bij aankomst op het doelscherm, wat de drempel voor de gebruiker aanzienlijk verlaagt.
	\item \textbf{Hiërarchische Taakstructuur via JSON:} Door de complexe onboarding via de metadata-structuur op te delen in hoofdtaken en substappen, blijft het proces overzichtelijk. Dit voorkomt dat de beheerder overweldigd raakt door de hoeveelheid configuraties en bevordert een stapsgewijze, foutloze inrichting van de applicatie.
	\item \textbf{Automatische Voortgangsdetectie:} In plaats van een oppervlakkige check op het bestaan van records, voert de Apex-backend \textit{field-level validation} uit. Het dashboard toont pas een succes-status wanneer de vereiste velden daadwerkelijk inhoudelijk correct zijn ingevuld, wat de betrouwbaarheid van de voortgangsrapportage garandeert.
	\item \textbf{Demo Data Generator:} Deze functionaliteit lost het 'empty state'-probleem op. Door de omgeving direct te vullen met coherente testdata via asynchrone Apex-processen, ervaart de gebruiker onmiddellijk de visuele waarde van componenten zoals de Team Cockpit (\textit{time-to-value}).
	\item \textbf{Cross-object Navigatie:} De implementatie ondersteunt processen die zich over meerdere Salesforce-apps en tabbladen uitspreiden. De walkthroughs blijven actief tijdens deze transities, waardoor de gebruiker nooit de draad van het proces verliest, ongeacht de locatie in het systeem.
\end{itemize}

\subsubsection*{Analyse van de 'Should Haves'}
De 'Should Haves' zijn grotendeels operationeel en richten zich specifiek op het verhogen van de gebruiksvriendelijkheid:

\begin{itemize}
	\item \textbf{Contextueel Linken van Documentatie:} Dankzij de koppeling via de \textit{Lightning Message Service} (LMS) fungeert de documentatie-viewer als een proactieve assistent. Het elimineert 'context switching'; de juiste handleiding volgt de acties op het dashboard synchroon, zonder dat de gebruiker de applicatie hoeft te verlaten.
	\item \textbf{Taakbevestiging (Manual Marking):} Voor stappen die buiten het bereik van automatische database-controles vallen—zoals specifieke instellingen in de Salesforce Setup-omgeving—biedt de handmatige taakbevestiging uitkomst. Door de status in de \textit{localStorage} te bewaren, behoudt de gebruiker de controle over het proces.
	\item \textbf{Data Reset-functie:} De mogelijkheid om de omgeving volledig op te schonen (\textit{Clear All Data}) creëert een veilige 'zandbak'-ervaring. Gebruikers kunnen vrij experimenteren met configuraties en testdata, wetende dat zij met één actie kunnen terugkeren naar een schone lei.
	\item \textbf{Custom Branding en Media:} Door strikt gebruik te maken van SLDS is de oplossing visueel consistent met de 24flow-huisstijl. De integratie van video-tutorials in de \textit{docked prompts} biedt bovendien een extra laag van ondersteuning voor de meest complexe configuratieonderdelen.
	\item \textbf{CSV-Upload (Placeholder):} De interface voor het opladen van eigen data is als functionele placeholder aanwezig. De UI-component (modal) is operationeel, maar de feitelijke data-mapping is wegens de scope van dit project gereserveerd voor toekomstige ontwikkeling.
\end{itemize}

\subsubsection*{Evaluatie van de 'Could Haves'}
Hoewel lager geprioriteerd, zijn \textbf{hervatten} en \textbf{telemetrie} volledig geïmplementeerd. De 'hervat'-logica zorgt ervoor dat een gebruiker nooit opnieuw hoeft te beginnen na een onderbreking. De telemetrie-logica (\texttt{logAction}) legt een fundament voor toekomstige analyses van het gebruikersgedrag, wat essentieel is voor de doorontwikkeling van het product.

\subsection{Validatie van de Use Cases}
De oplossing is getoetst aan de drie praktijkscenario's:
\begin{enumerate}
	\item \textbf{Initiële configuratie:} De testgebruiker doorliep de volledige setup van de Order Pool. De \textit{targeted prompts} zorgden ervoor dat de aandacht direct op de juiste configuratievelden lag.
	\item \textbf{Data Management:} Het genereren en wissen van data verliep vlekkeloos. De asynchrone afhandeling zorgde ervoor dat de gebruiker de interface kon blijven gebruiken terwijl de database-acties liepen.
	\item \textbf{Hervatten:} De status-logica identificeerde correct een afgebroken walkthrough en bood direct de 'Restart'-optie aan, wat de continuïteit van de onboarding waarborgde.
\end{enumerate}

\subsection{Technische Performantie en Beperkingen}
Het systeem opereert ruim binnen de Salesforce \textit{governor limits}. De JSON-gebaseerde architectuur garandeert dat de metadata-belasting minimaal blijft. 
\\
De voornaamste \textbf{beperking} blijft van technische aard: Salesforce staat niet toe dat In-App Guidance wordt weergegeven op administratieve 'Setup'-pagina's. Hoewel de documentatie-viewer en de directe navigatie dit deels opvangen, blijft dit een fundamentele platformrestrictie die een 100\% volledige in-app begeleiding verhindert.

%\begin{listing}
%  \begin{minted}{python}
%    import pandas as pd
%    import seaborn as sns
%
%    penguins = sns.load_dataset('penguins')
%    sns.relplot(data=penguins, x="flipper_length_mm", y="bill_length_mm", hue="species")
%  \end{minted}
%  \caption[Voorbeeld codefragment]{Voorbeeld van het invoegen van een codefragment.}
%\end{listing}




